\specialchapter{记号索引}

参考编号依次表示章节、段落和编号（或有时表示习题）。

\([x,y]\) \((x,y\) 李代数的元素): I.1.2.

\[
\mathfrak {g} ^ {0}
\]

\[
\mathfrak {g l} (\mathrm{E}), \mathfrak {g l} (n, \mathrm{K})
\]

\([\mathfrak{a},\mathfrak{b}],[z,\mathfrak{a}],[\mathfrak{a},z]\) (a, b 为子模, \(z\) 为李代数的元素): I.1.4.

af(M) (M 为 K-模): I.1.8.

\(U_{+}, U_{0}\) (U 为李代数的包络代数): I.2.1.

\(x_{\mathbb{M}}\) ( \(x\) 为李代数 \(\mathfrak{g}\) 的元素, M 为 \(\mathfrak{g}\) -模): I.3.1.

\[
\mathrm{T} _ {n}, \mathrm{U} _ {n}, \mathrm{G} ^ {n}: \text { I.2.6. }
\]

\(C(\rho)\) ( \(\rho\) 为李代数的表示): I.7.1.

\[
\mathcal {C} ^ {\infty} \mathfrak {g}, \mathcal {D} ^ {\infty} \mathfrak {g}: \mathrm{I}. 1, \text { Exercise 14. }
\]

\(x^{[p]}\) : I.1, 习题 20.

\(\mathbf{GL}(n,\mathbf{K})\) (形式群): I.1, 习题 25.

\(\mathfrak{o}(\Phi)\) : I.1, 习题 26.

\[
\mathrm{C} ^ {p} (\mathfrak {g}, \mathrm{M}), \mathrm{C} ^ {*} (\mathfrak {g}, \mathrm{M}), i (y), \theta (x), d, \mathrm{Z} ^ {p} (\mathfrak {g}, \mathrm{M}), \mathrm{B} ^ {p} (\mathfrak {g}, \mathrm{M}), \mathrm{H} ^ {p} (\mathfrak {g}, \mathrm{M}), \mathrm{H} ^ {*} (\mathfrak {g}, \mathrm{M})
\]

\(\mathfrak{sp}(2n,\mathbf{K})\) : I.6, 习题 25.

K: II. 约定.

g, U = Ug, σ: g → Ug: II.1.

E, E\(^{+}\): II.1.1. P(E), π, γ, c\(^{+}\): II.1.6. S(g), ε, γ: II.1.5. f\(_{E}\): U(P(E)) → E: II.1.6. M(X), l(w), Lib(X) = Lib\(_{K}\)(X): II.2.1. L(X) = L\(_{K}\)(X): II.2.2. φ: X → L(X): II.2.2. (a, r): II.2.3. L(u): II.2.5. Lib\(^{b}\)(X), L\(^{b}\)(X), L\(^{n}\)(X): II.2.6. P\(_{n}\): II.2.7. C\(^{n}\)(g): II.2.7. H, d\(_{y}\): II.2.10. w = Ψ(w): II.2.11. A(X) = A\(_{K}\)(X), A\(^{+}\)(X), Mo(X): II.3. π: II.3.2. (G\(_{\alpha}\)), (G\(_{\alpha}\) \(^{+}\)): II.4.1. v: II.4.2. gr(G), gr\(_{\alpha}\)(G): II.4.3. F(X), A(X), A\(^{n}\)(X): II.5. Â(X), ω: II.5.1. ε(a): II.5.2. l(x), exp(x), log(y): II.6.1. e(X), l(X): II.6.1. L(X): II.6.2. a ⊕ b: II.6.2. H, H\(_{n}\), H\(_{ra}\): II.6.4. H, Ω: II.7.2. A, exp\(_{A}\), log\(_{A}\), P(A\(^{I}\), A): II.7.3. v, θ = \(\frac{1}{p - 1}\): II.8. S(n): II.8.1. h(x, y): II.8.3. G\(_{R}\): II.8.4. μ(n): II.App. e, e\(_{G}\), γ(g), δ(g), Int(g), f: III.Conventions. GL(E), GL(n, K): III.1.1, III.3.10. G\(^{v}\): III.1.2. τ(g), ρ(x): III.1.5. (G, g, θ, m): III.1.10. T(m): III.2.1. T(G), T(φ): III.2.2. t * t': III.3.1, III.3.18.

U(G), U \(^{+}\) (G), U \(_{a}\) (G), U \(^{+}\) (G): III.3.1, III.3.18. T \(^{(y)}\) (G), T \(^{(\omega)}\) (G), T \(^{(\infty)}\) (G): III.3.1, III.3.18. t * f: III.3.4, III.3.18. D \(_{i}\) : III.3.5, III.3.18. L \(_{t}\) , R \(_{t}\) : III.3.6, III.3.18. L(G): III.3.7, III.3.18. L(φ): III.3.8, III.3.18. \textless{} t,f\textgreater{} : III.3.9, III.3.18. SL(E): III.3.10. Ad, Ad(g) : III.3.12. {[}α{]} \(^{2}\) : III.3.14. mod(ω) \(_{u}\) , mod φ: III.3.16. f \(^{-1}\) . df: III.3.17, III.3.18. H: III.4.2. g \(^{t}\) , ω(x) : III.4.3. x.y, x{[}n{]}: III.5. c \(_{\alpha\beta y}\) , B(x,y) : III.5.1. e \(_{\alpha}\) : III.5.2. ψ \(_{j}\) , ψ \(_{p,m}\) , \(\binom{l}{i}\) : III.5.3. E(x), L(x) : III.5.4. P \(_{M,x,y}\) : III.6.2. Ad(a) = Int(a): III.6.2. exp, exp \(_{G}\) : III.6.4. Ad(G) = Int(L(G)): III.6.4. L(ρ): III.6.5. G: III.6.10. A, m, p: III.7. G(a): III.7.4. h \(_{n}\) : III.7.5. G \(_{f}\) , log \(_{G}\) , log: III.7.6. D \(^{i}\) G, C \(^{i}\) G: III.9.1. Z \(_{G}\) (A), Z \(_{G}\) (a), λ \(_{g}\) (A), λ \(_{g}\) (a): III.9.3. N \(_{G}\) (A), N \(_{G}\) (a), n \(_{p}\) (a): III.9.4. R, N, r, n: III.9.7. π \(_{1}\) ⊗⋯⊗ π \(_{n}\) , T(π), S(π), ∧(π), T \(^{n}\) (π), S \(^{n}\) (π), ∧ \(^{n}\) (π): III.App.

\specialchapter{术语索引}

参考编号依次表示章节、段落和编号（或有时表示习题）

实或复李群的伴随群: III.6.4. 李代数中元素的线性映射: I.1.2. 李代数的表示: I.3.1. 李群的表示: III.3.12.

阿多定理: I.7.3.

代数（不一定结合）: I.1.1. 由代数扩张标量得到的代数: I.1.1. 李代数的包络代数: I.2.1. 李代数: I.1.2. 相反代数: I.1.1. 乘积代数: I.1.1. 商代数: I.1.1. 李 p-代数的受限包络代数: I.2, 习题 6. 模的对称代数: I.2.5.

几乎单李群: III.9.8.

n 次交错式: II.2.6.

李代数包络代数的主反自同构: I.2.4.

与 g-模（一个表示）相关的（双线性型）: I.3.6. 与运算律相关的（无穷小运算律）: III.3.7.

自由结合代数: II.3.

李代数的特殊自同构: I.6.8.

基本交换子: II.5.4. 李代数的族: II.2.3.

比伯巴赫定理: III.4, 习题 13.

双代数: II.1.2.

双不变截面: III.3.13.

二项式多项式: II.5, 习题 4.

括号: I.1.2.

左典范微分形式: III.3.13, III.3.18. 卡当判据: I.5.4. Casimir 元: I.3.7. 中心扩张: I.1.7. 群的滤过: II.4.4. 中心化子: I.1.6, III.9.3. 李代数的中心: I.1.6. 特征理想: I.1.4. 幂零类: II.2.7. 简单表示的类: I.3.1. 取值于 g-模的余边界、余链、余循环: I.3, 习题 12. 余代数: II.1.1. 交换李代数: I.1.3. 相容的群结构和流形结构: III.1.1. 完全不变双线性型: I.3.6. 可约表示: I.3.1. 复李群: III.1.1, III.8.1. 实李群的复化: III.6.10. g-模的同型分量: I.3.1. 半单李代数的简单分量: I.6.2. 复李群的共轭李群: III.1.1. 李代数包络代数元素的常数项: I.2.1. Magnus 代数元素的项: II.5.2. 代数相对于基的结构常数: I.1.1. 含有简单表示 n 次的表示: I.3.1. 解析表示的逆表示: III.3.11. 点分布与函数的卷积: III.3.4. 卷积乘积: III.3.1, III.3.18. 余单位: II.1.1. 李群的 C'-连通子集: III.6.2. 判据, 卡当: I.5.4.

代数的导子: I.1.1. 李代数的内导子: I.1.2. 相对于形式群的左不变导子: I.1, 习题 24. 自由群代数中的偏导子: II.5, 习题 2. （代数）通过扩张标量由代数导出: I.1.1. （表示）通过扩张标量由表示导出: I.3.8. 导出理想: I.1.5. 导出列: I.1.5.

映入李群的映射的左微分: III.3.17, III.3.18. 表示的维数: I.3.1. 对偶表示: I.3.3.

元素, Casimir: I.3.7. 幂零: I.6.3. 李代数包络代数所关联的分次代数中的 n 次元: I.2.6. 李代数包络代数中滤过 \(\leqslant n\) 的元: I.2.6. 半单: I.6.3. 李代数中的可交换元素: I.1.3. 消去定理: II.2.9. 恩格尔定理: I.4.2. 李代数的包络代数: I.2.1. 包络双代数: II.1.4. 李代数的等价扩张: I.1.7. 穷尽滤过: II.4.1. 指数: II.6.1, III.4.3, III.6.4. 扩张, 中心: I.1.7. 非本质: I.1.7. g-模的扩张, g 的表示的扩张: I.7.2. 李代数被李代数扩张: I.1.7. 平凡: I.1.7. 扩张, 等价: I.1.7.

忠实表示: I.3.1. 点分布的域: III.3.5, III.3.18. 滤过双代数: II.1.3. 群上的实滤过: II.4.1. 第一类的典范坐标图: III.4.3. 第一类的典范坐标系: III.4.3. 与李子代数相关的左叶层: III.4.1. 形式, 双线性, 与 g-模（g 的一个表示）相关: I.3.6. 双线性, 完全不变: I.3.6. 双线性, 不变: I.3.6. Killing: I.3.6. 公式, Jacobson: I.1, 习题 19. 自由李代数: II.2.2. 李 p-代数: II.3, 习题 4. 原群: II.2.1. 函数, 与滤过相关的阶: II.4.2.

与滤过群相关的分次群: II.4.3. 相关李代数: II.4.4. L(I) 的总分次: II.2.6. 群的交换形式: I.1, 习题 24. 形式群: I.1, 习题 24. 形式线性群: I.1, 习题 25. 形式正交群: I.1, 习题 26. 形式辛群: I.1, 习题 26. 李群芽: III.1.10. 由李代数定义的李群芽: III.4.2.

Hall 基: II.2.11. 公式: II.5, 习题 9. 集: II.2.10. Hausdorff 函数: II.7.2, II.8.3. 群: II.6.2. 逆元公式: II.6, 习题 4. 级数: II.6.4. 李代数的全自同构群: I.1, 习题 16. 同态, 代数: I.1.1. 从 g 的李子代数的包络代数到 g 的包络代数的典范同态: I.2.3. 从 g 的对称代数到与 g 的包络代数相关的分次代数上的典范满同态: I.2.6. 形式同态: I.1, 习题 24.

理想, 特征: I.1.4. 导出: I.1.5. 代数的左（右、双边）理想: I.1.1. 李代数的理想: I.1.4. 恒等式, Jacobi: I.1.2. 诱导的李群结构: III.4.5. 非本质扩张: I.1.7. 无穷小自同构: III.10.1. 内导子: I.1.2. 群的整滤过: II.4.1. 李群的子群: III.6.2. 不变双线性型: I.3.6. 左不变点分布域: III.3.6. g-模（一个表示）的不变: I.3.5. 不变截面: III.3.13.

逆像李群结构: III.1.9. 不可约表示: I.3.1. 同构表示: I.3.1. g 的对称代数到其包络代数底层向量空间上的典范同构: I.2.7. g-模的同型分量: I.3.1.

Jacobi 恒等式: I.1.2. Jordan 定理: III.4, 习题 11.

扩张的核: I.1.7. Killing 型: I.3.6.

g-模（表示）的最大幂零性理想: I.4.3. 李代数的最大幂零理想: I.4.4. 运算律片段: III.1.11. 形式群的运算律: I.1, 习题 24. 无穷小运算律: III.3.7, III.3.18. 自由原群元素的长度: II.2.1. Levi 子代数: I.6.8. Levi-Malcev 定理: I.6.8. 李代数: I.1.2. 交换李代数: I.1.3. 特征幂零李代数: I.4, 习题 19. 幂零李代数: I.4.1. 形式群的李代数: I.1, 习题 24. 李群的李代数: III.3.7. 李群芽的李代数: III.3.18. 模的端同态代数: I.1.2. 约化李代数: I.6.4. 半单李代数: I.6.1. 单李代数: I.6.2. 可解李代数: I.5.1. 李群: III.1.1. Lie 定理: I.5.3. 局部同构的群芽: III.1.10. 对数: II.6.1, III.7.6. 下中心列: I.1.5.

Magnus 代数: II.5.1. Magnus 群: II.5.2. 从李代数到其包络代数的典范映射: I.2.1. Maurer-Cartan 公式: III.3.14, III.3.18.

Möbius 函数: II, 附录. 反演公式: II, 附录. 态射, 李群: III.1.2. 李群芽: III.1.10. L(I) 的多重分次: II.2.6.

幂零李代数: I.4.1, II.2.7. 李代数的根基: I.5.3. 可赋范代数: II.7. 正规化子: I.1.4, III.9.4. 赋范李代数: II.7, II.8.2.

相反代数: I.1.1. 滤过群中元素的阶: II.4.2.

p-进李群: III.1.1, III.8.1. 李 p-代数: I.1, 习题 20. p-幂单李代数: I.4, 习题 23. 交换李 p-代数的 p-核: I.1, 习题 23. p-导子: I.2, 习题 7. 幂零群的 P-包络: II.4, 习题 15. p-同态: I.1, 习题 20. p-理想: I.1, 习题 22. P-整数: II.4, 习题 14. p-映射: I.1, 习题 20. p-多项式: I.7, 习题 5. 幂零群子群的 P-饱和: II.4, 习题 14. P-挠群, P-无挠群: II.4, 习题 14. 可交换元素: I.1.3, III.9.3. Poincaré-Birkhoff-Witt 定理: I.2.7. 李多项式: II.2.4. 多项式映射: II.2.4. 李代数的呈示: II.2.3. 双代数的本原元: II.1.2. 乘积代数: I.1.1. 李群的乘积: III.1.4. 李代数的半直积: I.1.8. 表示的张量积: I.3.2, III, 附录. 种类 N 的纯 g-模: I.3.1. 纯表示: I.3.1.

李拟子群: III.1.3.

商代数: I.1.1. 李群的商: III.1.6. 商表示: I.3.1.

李代数的幂零根基: I.5.3. 李代数的根基: I.5.2. 李群的根基: III.9.7.

实李群: III.1.1, III.8.1, III.8.2.

约化李代数: I.6.4. 李代数中的（李子代数）: I.6.6.

关系子: II.2.3.

端同态的复本: I.5, 习题 14.

表示, 伴随: I.3.1. 李群的解析线性表示: III.1.2. 完全可约: I.3.1. 含有简单表示 n 次: I.3.1. 对偶: I.3.3. 忠实: I.3.1. 不可约: I.3.1. 通过扩张标量环得到的: I.3.8. 李代数的: I.3.1. 种类 σ 的纯表示: I.3.1. 商: I.3.1. 半单: I.3.1. 简单: I.3.1.

表示, 同构: I.3.1. 相似: I.3.1.

标量限制, 由李群导出的李群: III.1.1.

可解李代数的根: III.9, 习题 17.

第二类的典范坐标图: III.4.3. 第二类的典范坐标系: III.4.3.

向量丛的截面: III.1.8.

李代数的半直积: I.1.8. 李群的半直积: III.1.4.

半单李代数: I.6.1. 半单李群: III.9.8. 半单表示: I.3.1.

分离滤过: II.4.1.

列, 连接两个子代数的合成列: I.1, 习题 14. 导出列: I.1.5. 李形式幂级数: II.6.3. 李代数的下中心列: I.1.5, II.2.7. 李代数的上中心列: I.1.6.

相似表示: I.3.1. 半单李代数的简单分量: I.6.2. 简单李代数: I.6.2. 简单表示: I.3.1. 可解李代数: I.5.1. 取值于 g-模的上同调空间: I.3, 习题 12. 李齐性空间: III.1.6. 特殊自同构: I.6.8. 纯 g-模的同型分量的种类: I.3.1. 标准群: III.7.3. 子代数: I.1.1. 李子代数: I.6.8. 李代数中的约化子代数: I.6.6. 次不变子代数: I.1, 习题 14. 李子群芽: III.1.10. 李子群: III.1.3. 子表示: I.3.1. 表示的直和: I.3.1. 模的对称代数: I.2.5.

t 次幂映射: III.4.3. 复合的切运算律: III.2.1. 李子群: III.4.5. 定理, 阿多: I.7.3. 恩格尔: I.4.2. Levi-Malcev: I.6.8. Lie: I.5.3. Poincaré-Birkhoff-Witt: I.2.7. Weyl: I.6.2. Zassenhaus: I.7.2. 传递子: I.1, 习题 1. 群, 严格上三角: II.4.6. 平凡扩张: I.1.7. g-模: I.3.1. 向量 G-丛: III.1.8. T(G) 的右（分别左）平凡化: III.2.1, III.2.2. 双边理想, I.1.1. 类型 (N) 的实李群: III.9, 习题 29.

\(u\) -余代数的 u-本原元: II.1.1.\\
幂单端同态: III.9.5.\\
连通李群的万有覆盖: III.1.9.

群的上中心列: II.4, 习题 18. 李代数的中心列: I.1.6.

Weyl 定理: I.6.2.

Zassenhaus 定理: I.7.2.

\summarychapter{总结 关于特征为 0 的域上的有限维李代数的某些性质}{总结}{关于特征为 0 的域上的有限维李代数的某些性质.}


设 g 是一个李代数, r 是其根基, n 是其最大幂零理想, s 是其幂零根基, t 是相对于 Killing 型的 g 的正交补. 于是 r, n, s, t 是特征理想, 且 \(r \supset t \supset n \supset s\) .

\begin{enumerate}
\def\labelenumi{(\Roman{enumi})}
\tightlist
\item
  以下性质中的任意一个都刻画半单李代数:
\end{enumerate}

\begin{enumerate}
\def\labelenumi{(\arabic{enumi})}
\tightlist
\item
  \(r = \{0\}\) ; (2) \(n = \{0\}\) ; (3) \(t = \{0\}\) ; (4) g 的每个交换理想都是零; (5) 代数 g 同构于简单李代数的一个乘积; (6) g 的每个有限维表示都是半单.
\end{enumerate}

\begin{enumerate}
\def\labelenumi{(\Roman{enumi})}
\setcounter{enumi}{1}
\tightlist
\item
  以下性质中的任意一个都刻画约化李代数:
\end{enumerate}

\begin{enumerate}
\def\labelenumi{(\arabic{enumi})}
\tightlist
\item
  \(\mathfrak{s} = \{0\}\) ; (2) \(\mathfrak{r}\) 是 \(\mathfrak{g}\) 的中心; (3) \(\mathcal{D}\mathfrak{g}\) 是半单; (4) \(\mathfrak{g}\) 是一个半单代数与一个交换代数的乘积; (5) \(\mathfrak{g}\) 的伴随表示是半单的; (6) \(\mathfrak{g}\) 有一个有限维表示, 使得相应的双线性型非退化; (7) \(\mathfrak{g}\) 有一个有限维半单忠实表示.
\end{enumerate}

\begin{enumerate}
\def\labelenumi{(\Roman{enumi})}
\setcounter{enumi}{2}
\tightlist
\item
  以下性质中的任意一个都刻画可解李代数:
\end{enumerate}

\begin{enumerate}
\def\labelenumi{(\arabic{enumi})}
\tightlist
\item
  \(\mathcal{D}^{p}\mathfrak{g} = \{0\}\) 当 p 充分大时; (2) 存在 g 的理想组成的递减序列 \(g = g_{0} \supset g_{1} \supset \cdots \supset g_{n} = \{0\}\), 使得代数 \(g_{i}/g_{i+1}\) 是交换的; (3) 存在一个递减序列
\end{enumerate}

\[
\mathfrak {g} = \mathfrak {g} _ {0} ^ {\prime} \supset \mathfrak {g} _ {1} ^ {\prime} \supset \dots \supset \mathfrak {g} _ {n ^ {\prime}} ^ {\prime} = \{0 \}
\]

\(\mathfrak{g}\) 的子代数, 使得 \(\mathfrak{g}_{i+1}^{\prime}\) 是 \(\mathfrak{g}_i^{\prime}\) 的理想, 且 \(\mathfrak{g}_i^{\prime}/\mathfrak{g}_{i+1}^{\prime}\) 是交换的; (4) 存在一个递减序列 \(\mathfrak{g} = \mathfrak{g}_0'' \supset \mathfrak{g}_1'' \supset \cdots \supset \mathfrak{g}_{n''}'' = \{0\}\) 的 \(\mathfrak{g}\) 的子代数, 使得 \(\mathfrak{g}_{i+1}''\) 是 \(\mathfrak{g}_i''\) 中余维 1 的理想; (5) \(\mathfrak{g} \supset \mathcal{D}\mathfrak{g}\) ; (6) \(\mathcal{D}\mathfrak{g}\) 是幂零的.

\begin{enumerate}
\def\labelenumi{(\Roman{enumi})}
\setcounter{enumi}{3}
\tightlist
\item
  以下性质中的任意一个都刻画幂零李代数:
\end{enumerate}

\begin{enumerate}
\def\labelenumi{(\arabic{enumi})}
\tightlist
\item
  \(\mathcal{C}_{g}^{p} = \{0\}\) 当 p 充分大时; (2) \(\mathcal{C}_{p}\mathfrak{g} = \mathfrak{g}\) 当 p 充分大时; (3) 存在 g 的理想组成的递减序列 \(g = g_{0} \supset g_{1} \supset \cdots \supset g_{p} = \{0\}\), 使得 \([g, g_{i}] \subset g_{i+1}\) ; (4) 存在 g 的理想组成的递减序列 \(g = g_{0}' \supset g_{1}' \supset \cdots \supset g_{p'}' = \{0\}\), 使得 \([g, g_{i}'] \subset g_{i+1}'\), 且 \(g_{i}'/g_{i+1}'\) 是 1-维的; (5) 存在一个整数 i, 使得
\end{enumerate}

\[
\left(\operatorname{ad} x _ {1}\right) \circ \left(\operatorname{ad} x _ {2}\right) \circ \dots \circ \left(\operatorname{ad} x _ {i}\right) = 0
\]

对所有 \(x_{1},\ldots ,x_{i}\), 它们都在 \(\mathfrak{g}\) 中 ; (6) 对所有 \(x\in \mathfrak{g}\) , ad \(x\) 都是幂零的.

\begin{enumerate}
\def\labelenumi{(\Alph{enumi})}
\setcounter{enumi}{21}
\tightlist
\item
  g 交换 \(\Rightarrow\) g 幂零 \(\Rightarrow\) g 的 Killing 型为零 \(\Rightarrow\) g 可解.
\end{enumerate}

g 交换 \(\Rightarrow\) g 约化.

g 半单 \(\Rightarrow\) g 约化.

\textbf{(VI) r 的刻画:}

\begin{enumerate}
\def\labelenumi{(\arabic{enumi})}
\tightlist
\item
  r 是 g 的最大可解理想; (2) r 是使 g/r 为半单的最小理想; (3) r 是使 g/r 为半单的唯一可解理想; (4) r 是相对于 Killing 型的 \(\mathcal{D}\mathfrak{g}\) 的正交补.
\end{enumerate}

\textbf{(VII) n 的刻画:}

\begin{enumerate}
\def\labelenumi{(\arabic{enumi})}
\tightlist
\item
  n 是 g 的最大幂零理想; (2) n 是 r 的最大幂零理想; (3) n 是 \(x \in r\) 中使 \(\operatorname{ad}_{\mathfrak{g}} x\) 为幂零的元素的集合; (4) n 是 \(x \in r\) 中使 \(\operatorname{ad}_{\mathfrak{r}} x\) 为幂零的元素的集合; (5) n 是 g 的最大理想, 使得对所有 \(x \in n\), \(\operatorname{ad}_{\mathfrak{g}} x\) 都是幂零的; (6) n 是 \(x \in g\) 中使 \(\operatorname{ad}_{\mathfrak{g}} x\) 属于由 1 和 \(\operatorname{ad}_{\mathfrak{g}} y\,(y \in g)\) 生成的结合代数的根基的元素的集合.
\end{enumerate}

\textbf{(VIII) s 的刻画:}

\begin{enumerate}
\def\labelenumi{(\arabic{enumi})}
\tightlist
\item
  \(\mathfrak{s}\) 是 \(\mathfrak{g}\) 的有限维简单表示的核的交; (2) \(\mathfrak{s}\) 是 \(\mathfrak{g}\) 的有限维半单表示的核中最小的一个; (3) \(\mathfrak{s}\) 是 \(\mathfrak{g}\) 的有限维表示的最大幂零性理想的交; (4) \(\mathfrak{s}\) 是 \(\mathfrak{g}\) 的最小理想, 使得 \(\mathfrak{g} / \mathfrak{s}\) 是约化的; (5) \(\mathfrak{s} = \mathfrak{r} \cap \mathcal{D}\mathfrak{g}\) ; (6) \(\mathfrak{s} = [\mathfrak{r}, \mathfrak{g}]\) ; (7) \(\mathfrak{s}\) 是相对于与 \(\mathfrak{g}\) 的有限维表示相关的双线性型的 \(\mathfrak{g}\) 的正交补的交.
\end{enumerate}


\specialchapter{尼古拉·布尔巴基 数学原理}

已出版的英文译本:

一般拓扑学, 第一部分

一般拓扑学, 第二部分

集合论

交换代数

代数, 第一部分

李群与李代数, 第一部分

\textbf{一般拓扑学, 第一部分}

\textbf{第一章. 拓扑结构}

§ 1. 开集, 邻域, 闭集

§ 2. 连续函数

§ 3. 子空间, 商空间

§ 4. 拓扑空间的乘积

§ 5. 开映射和闭映射

§ 6. 滤子

§ 7. 极限

§ 8. Hausdorff 空间和正则空间

§ 9. 紧空间和局部紧空间

§ 10. 真映射

§ 11. 连通性

习题

\textbf{第二章. 一致结构}

§ 1. 一致空间

§ 2. 一致连续函数

§ 3. 完备空间

§ 4. 一致空间与紧空间之间的关系

\textbf{第三章. 拓扑群}

§ 1. 群上的拓扑

§ 2. 子群, 商群, 同态, 齐性空间, 乘积群

§ 3. 群上的一致结构

§ 4. 在拓扑空间上适当作用的群; 拓扑群和带算子空间中的紧性

§ 5. 交换群中的无穷和

§ 6. 带算子的拓扑群; 拓扑环, 除环和域

\textbf{§ 7. 拓扑群和拓扑环的逆极限}

\textbf{第四章. 实数}

§ 1. 实数的定义

§ 2. 实直线的基本拓扑性质

§ 3. 实数域

§ 4. 扩展实直线

§ 5. 实值函数

§ 6. 连续和半连续实值函数

§ 7. 实数的无穷和与无穷积

§ 8. 实数的通常展开; R 的幂

习题

第五章. 单参数群 § 1. R 的子群和商群 § 2. 量的测量 § 3. 群 R 和 T 的拓扑刻画 § 4. 指数和对数 习题

第六章. 实数空间和射影空间 § 1. 实数空间 R \(^{n}\) § 2. 欧氏距离, 球和球面 § 3. 实射影空间 习题

第七章. 加法群 R \(^{n}\) § 1. R \(^{n}\) 的子群和商群 § 2. R \(^{n}\) 及其商群的连续同态 § 3. 群 R \(^{n}\) 中的无穷和 习题

第八章. 复数 § 1. 复数, 四元数 § 2. 角度测量, 三角函数 § 3. 复数的无穷和与无穷积 § 4. 复数空间和射影空间 习题

第九章. 实数在一般拓扑学中的应用 § 1. 由一族伪度量生成一致性; 可一致化空间 § 2. 度量空间和可度量化空间 § 3. 可度量化群, 赋值域, 赋范空间和代数 § 4. 正规空间 § 5. Baire 空间 § 6. Polish 空间, Souslin 空间, Borel 集 附录: 赋范代数中的无穷积 习题

第十章. 函数空间 § 1. S-收敛的一致性 § 2. 等度连续集 § 3. 特殊函数空间 § 4. 连续实值函数的逼近 习题

\textbf{一般拓扑学, 第二部分}

\textbf{集合论}

\textbf{第一章. 形式数学的描述}

§ 1. 项和关系

§ 2. 定理

§ 3. 逻辑理论

§ 4. 量化理论

§ 5. 等式理论

附录. 项和关系的刻画 习题

第二章. 集合论

§ 1. 集合化关系

§ 2. 有序对

§ 3. 对应

§ 4. 一族集合的并与交

§ 5. 一族集合的乘积

§ 6. 等价关系

习题

第三章. 有序集, 基数, 整数

§ 1. 序关系. 有序集

§ 2. 良序集

§ 3. 等势集. 基数

§ 4. 自然整数. 有限集

§ 5. 整数的性质

§ 6. 无限集

§ 7. 逆极限和直极限

习题

第四章. 结构

§ 1. 结构和同构

§ 2. 态射和导出结构

§ 3. 泛映射

习题

结果概要

引言

§ 1. 集合的元素和子集

§ 2. 函数

§ 3. 集合的乘积

§ 4. 一族集合的并、交、乘积

§ 5. 等价关系和商集

§ 6. 有序集

§ 7. 幂. 可数集

§ 8. 集合的尺度. 结构

\textbf{交换代数}

第一章. 平坦模 § 1. 图表和正合序列 § 2. 平坦模 § 3. 忠实平坦模 § 4. 平坦模与 ``Tor'' 函子 习题

第二章. 局部化 § 1. 素理想 § 2. 分式环和分式模 § 3. 局部环. 从局部到整体的过渡 § 4. 环的谱和模的支集 § 5. 有限生成投射模. 可逆分式理想 习题

第三章. 分次、滤过和拓扑 § 1. 有限生成分次代数 § 2. 滤过环和模的一般结果 § 3. 诺特环上的 m-进拓扑 § 4. 完备环中的提升 § 5. 滤过模的平坦性 习题

第四章. 关联素理想与初等分解 § 1. 与模相关的素理想 § 2. 初等分解 § 3. 分次模中的初等分解 习题

第五章. 整数 § 1. 整元的概念 § 2. 素理想的提升 § 3. 域上的有限生成代数 习题

第六章. 赋值 § 1. 赋值环 § 2. 位 § 3. 赋值 § 4. 赋值的高度 § 5. 由赋值定义的拓扑 § 6. 绝对值 § 7. 逼近定理 § 8. 赋值到代数扩张的延拓 § 9. 应用: 局部紧域 § 10. 赋值到超越扩张的延拓 习题

\textbf{交换代数}

第七章. 除子

§ 1. Krull 整环

§ 2. Dedekind 整环

§ 3. 唯一分解整环

§ 4. 整闭诺特整环上的模 习题

\textbf{代数, 第一部分}

\textbf{第一章. 代数结构}

§ 1. 复合律; 结合律; 交换律

§ 2. 单位元; 可消去元; 可逆元

§ 3. 作用

§ 4. 群和带算子的群

§ 5. 在集合上作用的群

§ 6. 扩张, 可解群, 幂零群

§ 7. 自由幺半群, 自由群

§ 8. 环

§ 9. 域

§ 10. 逆极限和直极限

习题

\textbf{第二章. 线性代数}

§ 1. 模

§ 2. 线性映射的模. 对偶性

§ 3. 张量积

§ 4. 张量积与同态模之间的关系

§ 5. 标量环的扩张

§ 6. 模的逆极限和直极限

§ 7. 向量空间

§ 8. 向量空间中标量域的限制

§ 9. 仿射空间和射影空间

§ 10. 矩阵

§ 11. 分次模和环

附录. 伪模

习题

\textbf{第三章. 张量代数、外代数、对称代数}

§ 1. 代数

§ 2. 代数的例子

§ 3. 分次代数

§ 4. 代数的张量积

§ 5. 张量代数. 张量

§ 6. 对称代数

§ 7. 外代数

§ 8. 行列式

§ 9. 范数和迹

§ 10. 导子

§ 11. 余代数, 多重线性型的乘积, 内积和对偶性 附录. 交代数. 八元数

习题
